\begin{examplebox}
	\textbf{\textcolor{CExample}{例1（基础）}}

	\textbf{\textcolor{CExample}{题目：}}
	在圆内接四边形 $ABCD$ 中，$\angle A = 80^\circ$，$\angle B = 110^\circ$，求 $\angle C$ 和 $\angle D$.

	\textbf{\textcolor{CExample}{解题步骤：}}
	\begin{description}[style=nextline, leftmargin=2.8em, labelsep=0.5em, itemsep=0.45em]
		\item[\textbf{第一步（利用互补）：}] 四边形 $ABCD$ 内接于圆，所以对角互补：
			\[
				\angle A + \angle C = 180^\circ,
				\qquad
				\angle B + \angle D = 180^\circ.
			\]
			\par\textbf{理由：}
			圆内接四边形对角互补结论.

		\item[\textbf{第二步（求 $\angle C$）：}] 将 $\angle A = 80^\circ$ 代入 $\angle A + \angle C = 180^\circ$，得
			\[
				\begin{aligned}
					\angle C &= 180^\circ - 80^\circ \\
					&= 100^\circ.
			\end{aligned}
			\]
			\par\textbf{理由：}
			移项计算.

		\item[\textbf{第三步（求 $\angle D$）：}] 将 $\angle B = 110^\circ$ 代入 $\angle B + \angle D = 180^\circ$，得
			\[
				\begin{aligned}
					\angle D &= 180^\circ - 110^\circ \\
					&= 70^\circ.
			\end{aligned}
			\]
			\par\textbf{理由：}
			移项计算.
	\end{description}

	\textbf{\textcolor{CExample}{关键结论：}}
	$\angle C = 100^\circ$，$\angle D = 70^\circ$.

	\medskip
	\textbf{\textcolor{CExample}{例2（稍变形）}}

	\textbf{\textcolor{CExample}{题目：}}
	在圆内接四边形 $ABCD$ 中，$\angle A = 2\angle C$，求 $\angle A$ 和 $\angle C$ 的度数.

	\textbf{\textcolor{CExample}{解题步骤：}}
	\begin{description}[style=nextline, leftmargin=2.8em, labelsep=0.5em, itemsep=0.45em]
		\item[\textbf{第一步（设未知数）：}] 设 $\angle C = x$，则 $\angle A = 2x$.
			\par\textbf{理由：}
			利用题目条件建立关系.

		\item[\textbf{第二步（应用互补）：}] 圆内接四边形对角互补，$\angle A + \angle C = 180^\circ$，代入得
			\[
				2x + x = 180^\circ.
			\]
			\par\textbf{理由：}
			对角互补结论.

		\item[\textbf{第三步（解方程）：}] 解得 $x = 60^\circ$，故 $\angle C = 60^\circ$，$\angle A = 120^\circ$.
			\[
				3x = 180^\circ,
				\qquad
				x = 60^\circ.
			\]
			\par\textbf{理由：}
			代数求解.
	\end{description}

	\textbf{\textcolor{CExample}{关键结论：}}
	$\angle A = 120^\circ$，$\angle C = 60^\circ$.

	\medskip
	\textbf{\textcolor{CExample}{例3（综合应用）}}

	\textbf{\textcolor{CExample}{题目：}}
	如图，圆内接四边形 $ABCD$ 中，延长 $BC$ 至 $E$，$\angle DCE = 70^\circ$，求 $\angle A$ 的度数.

	\textbf{\textcolor{CExample}{解题步骤：}}
	\begin{description}[style=nextline, leftmargin=2.8em, labelsep=0.5em, itemsep=0.45em]
		\item[\textbf{第一步（识图）：}] 由题意知 $\angle DCE$ 是四边形 $ABCD$ 的外角.
			\par\textbf{理由：}
			延长一边，得到外角.

		\item[\textbf{第二步（应用外角定理）：}] 圆内接四边形的外角等于其内对角，所以 $\angle DCE = \angle A$.
			\[
				\angle DCE = \angle A.
			\]
			\par\textbf{理由：}
			圆内接四边形外角定理（可由对角互补推出）.

		\item[\textbf{第三步（代值）：}] $\angle DCE = 70^\circ$，所以 $\angle A = 70^\circ$.
			\[
				\angle A = 70^\circ.
			\]
			\par\textbf{理由：}
			等量代换.
	\end{description}

	\textbf{\textcolor{CExample}{关键结论：}}
	$\angle A = 70^\circ$.
\end{examplebox}
